% !TEX root = ../main.tex

\documentclass[../main.tex]{subfiles}

\graphicspath{{\subfix{../images/}}}

\begin{document}
	\thispagestyle{empty}
	
	\noindent
	На №№ с., \TotalValue{totalfigures} рисунков, \TotalValue{totaltables} таблиц, №№ приложений.
	
	КЛЮЧЕВЫЕ СЛОВА: КРОССПЛАТФОРМЕННОЕ ПРИЛОЖЕНИЕ, KOTLIN MULTIPLATFORM, COMPOSE MULTIPLATFORM, ГЕЙМИФИКАЦИЯ, ТРЕКИНГ ДОСТИЖЕНИЙ, MVVM, KTOR, OPENAPI GENERATOR.
	
	Тема выпускной квалификационной работы: <<Разработка кроссплатформенного мобильного приложения для трекинга достижений с элементами геймификации>>. Работа посвящена созданию приложения Achi, позволяющего пользователям структурировать цели в виде паков достижений, отслеживать пошаговый прогресс и обмениваться наборами с другими пользователями.
	Задачи, которые решались в ходе исследования:
	\begin{enumerate}
		\item Провести анализ существующих решений в области трекинга достижений и выявить их функциональные и технические недостатки.
		\item Сформулировать функциональные и нефункциональные требования к разрабатываемому приложению.
		\item Обосновать выбор технологического стека (Kotlin Multiplatform, Compose Multiplatform, Ktor, Navigation~3, OpenAPI Generator) и спроектировать архитектуру на основе паттерна MVVM с Repository.
		\item Реализовать кроссплатформенное приложение с поддержкой иерархической структуры данных (паки~$\rightarrow$~достижения~$\rightarrow$~шаги), аутентификации, синхронизации прогресса с сервером и обмена паками по уникальным кодам.
		\item Провести функциональное и кроссплатформенное тестирование разработанного решения.
	\end{enumerate}
	
	Разработка выполнена в среде Android Studio с использованием Kotlin Multiplatform и Compose Multiplatform. Сетевой слой сгенерирован из OpenAPI-спецификации, внедрение зависимостей организовано через ручной DI-контейнер \texttt{AppModule}, навигация реализована с помощью Navigation~3. Приложение развёрнуто на платформах Android, iOS и Web (WasmJS).
	
	По результатам работы создано кроссплатформенное приложение, удовлетворяющее сформулированным требованиям: пользователь может управлять коллекцией паков достижений, отслеживать прогресс выполнения шагов с синхронизацией на сервер, создавать собственные паки и добавлять чужие по коду. Полученное решение демонстрирует практическую применимость современного KMP-стека для разработки продуктов в области персональной продуктивности с элементами геймификации.
	
\end{document}
